% Section 4: Evaluating Virtual Staining Models in Bioimaging

\section{Virtual Staining Evaluation Metrics}

\begin{frame}[fragile]{Pixel-Wise Metrics: PSNR \& MAE}
\begin{columns}
\column{0.48\textwidth}
\textbf{Mathematical Formulation}

Mean Absolute Error (MAE):
\begin{equation*}
\text{MAE}(x, \hat{x}) = \frac{1}{N} \sum_{i=1}^N |x_i - \hat{x}_i|
\end{equation*}

Peak Signal-to-Noise Ratio (PSNR in dB):
\begin{equation*}
\text{PSNR}(x, \hat{x}) = 10 \cdot \log_{10} \left( \frac{\text{MAX}_x^2}{\text{MSE}(x, \hat{x})} \right)
\end{equation*}

where $\text{MSE}(x, \hat{x}) = \frac{1}{N} \sum_{i=1}^N (x_i - \hat{x}_i)^2$.

\vspace{1mm}
\textbf{\color{commentgreen}Biological Intuition (DAPI $\to$ GFP):}\\
\fontsize{6.5pt}{7.5pt}\selectfont
MAE measures average GFP intensity error per pixel. Lower MAE means correctly predicting 0 GFP in negative cells and accurate brightness in GFP-positive cells.

\column{0.52\textwidth}
\textbf{Python Code Equivalent}
\begin{lstlisting}
def compute_psnr_mae(x_gt, x_pred, max_val=1.0):
    # MAE Calculation
    mae = torch.mean(torch.abs(x_gt - x_pred))
    
    # MSE & PSNR Calculation
    mse = torch.mean((x_gt - x_pred) ** 2)
    if mse == 0:
        psnr = torch.tensor(float('inf'))
    else:
        psnr = 10.0 * torch.log10((max_val ** 2) / mse)
        
    return psnr.item(), mae.item()
\end{lstlisting}
\end{columns}
\end{frame}


\begin{frame}[fragile]{Structural Similarity Index (SSIM)}
\begin{columns}
\column{0.48\textwidth}
\textbf{Mathematical Formulation}

SSIM models human visual perception (luminance, contrast, structure):
\begin{equation*}
\text{SSIM}(x, \hat{x}) = \frac{(2\mu_x \mu_{\hat{x}} + C_1)(2\sigma_{x\hat{x}} + C_2)}{(\mu_x^2 + \mu_{\hat{x}}^2 + C_1)(\sigma_x^2 + \sigma_{\hat{x}}^2 + C_2)}
\end{equation*}

\begin{itemize}
    \item $\mu_x, \mu_{\hat{x}}$: Local mean intensities.
    \item $\sigma_x^2, \sigma_{\hat{x}}^2$: Local variances.
    \item $\sigma_{x\hat{x}}$: Local covariance.
    \item $C_1, C_2$: Stabilization constants.
\end{itemize}

\vspace{1mm}
\textbf{\color{commentgreen}Biological Intuition (DAPI $\to$ GFP):}\\
\fontsize{6.5pt}{7.5pt}\selectfont
Evaluates whether predicted GFP spatial boundaries match true cell shapes around DAPI-stained nuclei, rather than just checking overall background brightness.

\column{0.52\textwidth}
\textbf{Python Code Equivalent}
\begin{lstlisting}
from skimage.metrics import structural_similarity as ssim

def compute_ssim_numpy(x_gt_np, x_pred_np):
    # x_gt_np, x_pred_np: [H, W] float32 arrays in [0, 1]
    ssim_score = ssim(
        x_gt_np, 
        x_pred_np, 
        data_range=1.0,
        win_size=7
    )
    return ssim_score
\end{lstlisting}
\end{columns}
\end{frame}


\begin{frame}[fragile]{Pearson Correlation Coefficient (PCC)}
\begin{columns}
\column{0.48\textwidth}
\textbf{Mathematical Formulation}

Linear correlation between predicted fluorescence $\hat{x}$ and ground truth $x$:
\begin{equation*}
\text{PCC}(x, \hat{x}) = \frac{\sum_{i=1}^N (x_i - \bar{x})(\hat{x}_i - \bar{\hat{x}})}{\sqrt{\sum_{i=1}^N (x_i - \bar{x})^2} \sqrt{\sum_{i=1}^N (\hat{x}_i - \bar{\hat{x}})^2}}
\end{equation*}

\begin{itemize}
    \item Range: $[-1, 1]$.
    \item $\text{PCC} = 1$: Perfect linear relationship.
    \item Insensitive to global linear intensity scaling/shifts!
\end{itemize}

\vspace{1mm}
\textbf{\color{commentgreen}Biological Intuition (DAPI $\to$ GFP):}\\
\fontsize{6.5pt}{7.5pt}\selectfont
Checks if relative GFP intensity rankings across cells match truth (strongly positive vs weakly positive cells) regardless of overall brightness scaling.

\column{0.52\textwidth}
\textbf{Python Code Equivalent}
\begin{lstlisting}
def compute_pcc(x_gt, x_pred):
    # Flatten spatial dimensions
    x1 = x_gt.flatten()
    x2 = x_pred.flatten()
    
    # Zero-mean vectors
    x1_m = x1 - torch.mean(x1)
    x2_m = x2 - torch.mean(x2)
    
    # Covariance divided by product of standard deviations
    num = torch.sum(x1_m * x2_m)
    den = torch.sqrt(torch.sum(x1_m ** 2) * torch.sum(x2_m ** 2) + 1e-8)
    
    pcc = num / den
    return pcc.item()
\end{lstlisting}
\end{columns}
\end{frame}


\begin{frame}[fragile]{Biological Validity: Single-Cell Infection Quantification}
\begin{columns}
\column{0.48\textwidth}
\textbf{VIRVS Biological Metric}

Virtual staining models must preserve biologically meaningful signal quantification!

For single-cell binary mask $M_k \in \{0, 1\}^{H \times W}$:
\begin{equation*}
I_{\text{cell}, k} = \frac{1}{\sum_{ij} M_{k, ij}} \sum_{i,j} M_{k, ij} \cdot x_{ij}
\end{equation*}

Benchmark compares true single-cell viral intensity $I_{\text{true}, k}$ vs predicted $I_{\text{pred}, k}$ across infection states.

\vspace{1mm}
\textbf{\color{commentgreen}Biological Intuition (DAPI $\to$ GFP):}\\
\fontsize{6.5pt}{7.5pt}\selectfont
Directly quantifies mean GFP reporter signal within each DAPI-segmented cell mask to classify individual cells as infected (GFP+) or uninfected (GFP-).

\column{0.52\textwidth}
\textbf{Python Code Equivalent}
\begin{lstlisting}
def compute_cell_level_viral_reporter(x_fluo, cell_masks):
    # cell_masks: [Num_Cells, H, W] boolean array
    cell_intensities = []
    
    for k in range(cell_masks.shape[0]):
        mask = cell_masks[k]
        # Mean fluorescence intensity inside cell mask k
        cell_intensity = (x_fluo * mask).sum() / (mask.sum() + 1e-8)
        cell_intensities.append(cell_intensity)
        
    return torch.tensor(cell_intensities)
\end{lstlisting}
\end{columns}
\end{frame}
